\chapter{\(\mathrm{L}_{\mathrm{ACD}}\) Subconditions: Archimedean Stratum, Colimit Commutation, Cross-Prime Hecke, Sph-Finite Preservation}
\label{v4-arith:lacd-sub}

\emph{The modular decomposition of \(\mathrm{L}_{\mathrm{ACD}}\)
(\Cref{v4-arith:acd:thm:modular-decomposition},
\Cref{v4-arith:sw:thm:LACD-split}) produces five subconditions:
\(\mathrm{L}_{\mathrm{ACD}}.0\) (archimedean stratum admissibility,
isolated by the archimedean-factor analysis of
\Cref{v4-arith:sw:wave2}),
\(\mathrm{L}_{\mathrm{ACD}}.1\) (dg-categorical colimit commutation),
\(\mathrm{L}_{\mathrm{ACD}}.2\) (cross-prime Hecke coherence
including \(U(\mathfrak{gl}_{2})^{K_{\infty}}\)-equivariance),
\(\mathrm{L}_{\mathrm{ACD}}.3\) (sph-finite preservation), and
\(\mathrm{L}_{\mathrm{ACD}}.4\) (strong approximation for
\(SL_{2}/\mathbb{Q}\) together with the determinant idele-class
quotient for \(GL_{2}/\mathbb{Q}\)). Kneser 1966 and
Platonov--Rapinchuk 1994 supply the simply connected part; Tate's
idele formalism supplies the determinant line. This chapter separates the formal
subconditions from the remaining comparison data: Harish-Chandra 1966
admissibility, Jacquet--Langlands SLN~114 archimedean theory, Vogan
1981 classification, Wallach 1988/1992 real-reductive-groups spectrum
theory, Bernstein--Zelevinsky 1976 induced representations, and
Lurie HA \S5.1, \S4.8 and Fresse 2017 \(E_{n}\) operadic machinery;
the Fargues--Scholze list is not invoked for the formal steps.
The \(\infty\)-categorical Bernstein lift and the Bernstein-centre
diamond-gluing compatibility remain open inputs.}

\section{Subconditions and Independent Sources}
\label{v4-arith:lacd:context}

\subsection{The Subcondition List}

The modular decomposition (\Cref{v4-arith:acd:thm:modular-decomposition})
splits \(\mathrm{L}_{\mathrm{ACD}}\) into \((L_{1})\)--\((L_{4})\);
\Cref{v4-arith:sw:thm:LACD-split} relabels these as
\(\mathrm{L}_{\mathrm{ACD}}.1\)--\(\mathrm{L}_{\mathrm{ACD}}.4\):
\begin{itemize}
\item \(\mathrm{L}_{\mathrm{ACD}}.1\) (colimit commutation);
\item \(\mathrm{L}_{\mathrm{ACD}}.2\) (cross-prime Hecke coherence);
\item \(\mathrm{L}_{\mathrm{ACD}}.3\) (sph-finite preservation);
\item \(\mathrm{L}_{\mathrm{ACD}}.4\) (strong approximation for
\(SL_{2}\) plus determinant idele-class quotient for \(GL_{2}\)).
\end{itemize}
The archimedean-factor analysis of \Cref{v4-arith:sw:wave2} identifies
a fifth subcondition \(\mathrm{L}_{\mathrm{ACD}}.0\) (archimedean
stratum admissibility) and tightens \(\mathrm{L}_{\mathrm{ACD}}.2\) to
include \(U(\mathfrak{gl}_{2})^{K_{\infty}}\)-equivariance.

The local list of subconditions considered here is:
\begin{itemize}
\item[\(\mathrm{L}_{\mathrm{ACD}}.0\):] archimedean stratum with
Harish-Chandra admissibility replacing Iwahori-fixedness;
\item[\(\mathrm{L}_{\mathrm{ACD}}.1\):] colimit commutes with
sph-finite truncation in \(\mathrm{Pr}^{\mathrm{L,st}}\);
\item[\(\mathrm{L}_{\mathrm{ACD}}.2\):] cross-prime Hecke coherence
including \(U(\mathfrak{gl}_{2})^{K_{\infty}}\)-equivariance;
\item[\(\mathrm{L}_{\mathrm{ACD}}.3\):] sph-finite preservation under
restricted-tensor colimit;
\item[\(\mathrm{L}_{\mathrm{ACD}}.4\):] strong approximation for
\(SL_{2}/\mathbb{Q}\), with the determinant quotient retained for
\(GL_{2}/\mathbb{Q}\).
\end{itemize}

\subsection{Independent Sources}

The formal parts of \(\mathrm{L}_{\mathrm{ACD}}.0\)--\(\mathrm{L}_{\mathrm{ACD}}.3\)
are treated here via an independent list; the Fargues--Scholze
geometrization construction is not invoked for those formal steps.

Primary sources used:
\begin{description}
\item[Harish-Chandra 1966] ``Discrete series for semisimple Lie
groups II'', Acta Math.\ 116. Admissible representations of
reductive Lie groups; the category of \((\mathfrak{g}, K)\)-modules.
\item[Jacquet--Langlands 1970] \emph{Automorphic Forms on \(GL(2)\)},
SLN 114. Archimedean local theory for \(GL_{2}(\mathbb{R})\) and
\(GL_{2}(\mathbb{C})\).
\item[Vogan 1981] \emph{Representations of Real Reductive Lie
Groups}, change in Math.\ 15. Langlands--Vogan classification at
\(\mathbb{R}\).
\item[Wallach 1988/1992] \emph{Real Reductive Groups I, II}, Pure
\& Applied Math., Academic Press. Admissible spectrum theory.
\item[Bernstein--Zelevinsky 1976] ``Induced representations of
reductive \(\mathfrak{p}\)-adic groups I'', Ann.\ Sci.\ ENS.
Bernstein decomposition of \(p\)-adic Hecke categories.
\item[Lurie HA] \emph{Higher Algebra}, 2017 draft. \S5.1 (\(E_{n}\)
monoidal coherence), \S4.8 (colimits in \(\mathrm{Pr}^{\mathrm{L}}\)).
\item[Fresse 2017] \emph{Homotopy of Operads and Grothendieck--Teichm\"uller
Groups, Vol.\ II}. Integral \(E_{n}\) operadic
machinery.
\end{description}

Secondary sources: Beilinson--Bernstein 1993 (localization;
D-modules on \(G/B\), archimedean localization), Nadler 2009
(microlocal sheaves on real flag varieties), Borel 1976
(admissible representations of semi-simple groups over local
fields), Kneser 1966 (strong approximation for simply-connected
groups), Platonov--Rapinchuk 1994 \emph{Algebraic Groups and
Number Theory}, Moeglin--Waldspurger 2016 (twisted trace
formula).

The initial arithmetic comparison derivation invoked Raskin 2021 JAMS for \((L_{1})\), a
Fargues--Scholze local input together with a function-field analogy for
\((L_{2})\), and Gaitsgory--Rozenblyum for
\(\mathrm{IndCoh}_{\mathrm{Nilp}}\) continuity. The comparison sources
here are Lurie HA \S5.1 +
\S4.8 (colimit commutation via pointed-\(E_{n}\) closure),
Harish-Chandra + Vogan + Wallach (archimedean admissibility), and
Bernstein--Zelevinsky + Lurie HA \S5.1.2 (symmetric-monoidal Hecke
coherence). The primary-source independence comparison passes for \(\mathrm{L}_{\mathrm{ACD}}.0\),
\(\mathrm{L}_{\mathrm{ACD}}.1\), and \(\mathrm{L}_{\mathrm{ACD}}.3\): no
primary author in the construction-2 list appears in the construction-1 list for
those subconditions; Raskin 2021 does not treat archimedean admissibility,
and Wallach 1988 does not construct \(\infty\)-categorical colimits in
\(\mathrm{Pr}^{\mathrm{L,st}}\). The \(\infty\)-categorical lift in the
Bernstein--Zelevinsky compatibility proposition is not used as a
disjoint verification. It is isolated below as a separate lift datum.

\section{Overview of Subcondition Closures}
\label{v4-arith:lacd:table}

The formal subconditions \(\mathrm{L}_{\mathrm{ACD}}.0\)--\(\mathrm{L}_{\mathrm{ACD}}.3\)
are analysed in the four sections that follow; the compatibility
with the completed Tate character
\(\Lambda_{\zeta}(s) = \Gamma_{\mathbb{R}}(s)\zeta(s)\)
functional equation occupies \Cref{v4-arith:lacd:xi-compat}.

\subsection{\(\mathrm{L}_{\mathrm{ACD}}.0\): archimedean stratum}

At the archimedean place, the local category
\(\mathrm{D}^{\mathrm{sph\text{-}fin}}([GL_{2}(\mathbb{R})])\)
is not defined by Iwahori-fixed vectors (as at finite primes) but by
\((\mathfrak{g}, K_{\infty})\)-admissibility: the empty-\(S\) base case
of the restricted tensor product is \(\mathrm{Rep}(K_{\infty})\), not Vec,
and any construction of the restricted tensor product
\(\bigotimes'_{v \in |\overline{C}|}\mathcal{C}_{v}\) must carry a
parallel archimedean factor.

The resolution is Harish-Chandra admissibility. The archimedean local
category is \(\mathrm{D}((\mathfrak{g}, K_{\infty})\text{-mod}^{\mathrm{sph\text{-}fin}})\)
with admissibility replacing Iwahori-fixedness.
Jacquet--Langlands SLN~114 \S1--3 and Vogan 1981 classify the admissible
\((\mathfrak{g}, K_{\infty})\)-modules for \(GL_{2}(\mathbb{R})\); the
sph-fin cut selects the modules finite over the archimedean Hecke algebra
(the convolution algebra of bi-\(K_{\infty}\)-finite distributions).
The detailed construction and proof appear in \Cref{v4-arith:lacd:L0}.

\subsection{\(\mathrm{L}_{\mathrm{ACD}}.1\): colimit commutation}

The filtered colimit defining \(\bigotimes'_{v}\mathcal{C}_{v}\) and
the sph-fin subcategory inclusion
\(\mathcal{C}_{v}^{\mathrm{sph\text{-}fin}} \hookrightarrow \mathcal{C}_{v}\)
are a priori unrelated operations; it is not immediate that they commute.

F2.c tightened (\Cref{v4-arith:tate:thm:F2c-main}) supplies the
universal property needed: the restricted tensor product of pointed
\(E_{n}\)-monoidal categories carries a canonical \(E_{n}\)-monoidal
structure, unique up to a contractible space of choices, compatible with
the insertion functors \(\iota_{S}\). Sph-finite truncation is defined by
compact generation (\Cref{v4-arith:acd:lem:sph-fin-preservation}), and
Lurie HA~5.3.2.9 gives that filtered colimits of compactly generated
categories along compact-generator-preserving functors are compactly
generated. Hence the colimit commutes with the sph-fin inclusion. The
proof appears in \Cref{v4-arith:lacd:L1}.

\subsection{\(\mathrm{L}_{\mathrm{ACD}}.2\): cross-prime Hecke coherence}

For distinct primes \(p \neq q\), the local Hecke operators \(T_{p}\)
and \(T_{q}\) act on disjoint tensor factors of the restricted tensor
product. The question is whether they commute up to coherent associators
and unit coherences, and whether this extends to
\(U(\mathfrak{gl}_{2})^{K_{\infty}}\)-equivariance at the archimedean place.

Lurie HA~5.1.2.4 (Dunn additivity) and HA~5.1.2.26 (centrality of the
unit in an \(E_{n}\)-monoidal category) supply the coherence: endofunctors
of a tensor product acting on disjoint tensor factors commute up to a
contractible space of coherence data. Extension to the archimedean place
is immediate from the disjointness of \(\mathcal{C}_{\infty}\) and each
\(\mathcal{C}_{p}\) as tensor factors. The proof appears in
\Cref{v4-arith:lacd:L2}.

\subsection{\(\mathrm{L}_{\mathrm{ACD}}.3\): sph-finite preservation}

Sph-finite truncation composes with the colimit to give the sph-finite
truncation of the colimit (not merely an approximation), because this is
the content of \(\mathrm{L}_{\mathrm{ACD}}.1\): if the colimit commutes with
sph-fin, then the sph-fin subcategory of the colimit equals the colimit of
sph-fin subcategories. \Cref{v4-arith:acd:lem:sph-fin-preservation}
handles the compact generation step. The proof appears in
\Cref{v4-arith:lacd:L3}.

\subsection{Archimedean character and the completed \(\xi\)-function}

The archimedean factor \(\Gamma_{\mathbb{R}}(s) = \pi^{-s/2}\Gamma(s/2)\)
is the archimedean local \(L\)-factor of the trivial representation;
Tate 1950 \S4.3 computes it as the Mellin transform of the Gaussian.
Thus the trace comparison, if it is constructed, must send the
archimedean unit to the Gaussian Tate vector. The restricted tensor
product is compared with
\(\Lambda_{\zeta}(s)=\Gamma_{\mathbb R}(s)\zeta(s)\), not with
\(\zeta(s)\) alone. The necessary Tate calculation appears in
\Cref{v4-arith:lacd:xi-compat}; the categorical trace comparison is a
separate input.

\section{\(\mathrm{L}_{\mathrm{ACD}}.0\): Archimedean Stratum Closure}
\label{v4-arith:lacd:L0}

\subsection{Statement and Construction}

\begin{definition}[archimedean admissible \((\mathfrak{g}, K_{\infty})\)-module]
\label{v4-arith:lacd:def:arch-admissible}
Let \(\mathfrak{g} = \mathfrak{gl}_{2}(\mathbb{R})\) and \(K_{\infty} = O(2)
\subseteq GL_{2}(\mathbb{R})\). A \emph{\((\mathfrak{g}, K_{\infty})\)-module}
is a real vector space \(V\) with a \(\mathfrak{g}\)-module structure and
a compatible \(K_{\infty}\)-action such that: (i) the \(K_{\infty}\)-action
is locally finite; (ii) the differentiated \(K_{\infty}\)-action agrees
with the restriction of the \(\mathfrak{g}\)-action to
\(\mathrm{Lie}(K_{\infty}) = \mathfrak{k}\). The module is \emph{admissible}
if also: (iii) each \(K_{\infty}\)-isotype \(V(\sigma)\) for
\(\sigma \in \widehat{K_{\infty}}\) is finite-dimensional (Harish-Chandra
1966 \S3).
\end{definition}

\begin{definition}[archimedean sph-finite category]
\label{v4-arith:lacd:def:arch-sph-fin}
The archimedean sph-finite category is
\begin{equation}
\mathcal{C}_{\infty} \;:=\; \mathrm{D}^{\mathrm{sph\text{-}fin}}
\bigl((\mathfrak{g}, K_{\infty})\text{-mod}^{\mathrm{adm}}\bigr)
\end{equation}
the stable presentable \(\infty\)-category of admissible
\((\mathfrak{g}, K_{\infty})\)-modules of finite type over the
archimedean Hecke algebra
\(\mathcal{H}_{\infty} = \mathrm{End}_{(\mathfrak{g}, K_{\infty})}
(V_{K_{\infty}}^{\mathrm{fin}})\), where \(V_{K_{\infty}}^{\mathrm{fin}}\)
ranges over admissible \((\mathfrak{g}, K_{\infty})\)-modules.
\end{definition}

\begin{theorem}[\(\mathrm{L}_{\mathrm{ACD}}.0\) closed]
\label{v4-arith:lacd:thm:L0-closed}
\ClaimStatusProvedHere. The archimedean category \(\mathcal{C}_{\infty}\)
of \Cref{v4-arith:lacd:def:arch-sph-fin} is:
\begin{enumerate}
\item a well-defined stable presentable \(\infty\)-category;
\item equipped with a distinguished unit
\(u_{\infty} \in \mathcal{C}_{\infty}\) (the \(K_{\infty}\)-spherical
unramified principal series) whose unit-comparison map
\(\eta_{\infty}: \mathbb{1}_{\mathcal{C}_{\infty}} \to u_{\infty}\)
is canonical;
\item classifiable by Langlands--Vogan parameters (Vogan 1981
Chapters 2--5): the isomorphism classes of irreducible admissible
\((\mathfrak{g}, K_{\infty})\)-modules correspond bijectively to
triples \((\lambda, \mu, \varepsilon)\) with \(\lambda \in \mathbb{C}\)
(weight), \(\mu \in \mathbb{Z}_{\geq 0}\) (discrete-series parameter),
and \(\varepsilon \in \{\pm 1\}\) (central character sign);
\item the empty-\(S\) base case of the restricted tensor product is
\(\mathrm{Rep}(K_{\infty})\), not Vec.
\end{enumerate}
\end{theorem}

\begin{proof}
\textbf{Step 1: Stability and presentability.} The category of
\((\mathfrak{g}, K_{\infty})\)-modules is abelian (Wallach 1988
\S0.1). Passing to the derived category and restricting to admissible
modules (finite per \(K_{\infty}\)-isotype) preserves stability (the
derived category of an abelian category is stable). Presentability
follows from Wallach 1988 Theorem 0.A.2: the category of admissible
\((\mathfrak{g}, K_{\infty})\)-modules is an abelian Grothendieck
category, so its derived category is presentable (Lurie HA 1.3.5.21).

\textbf{Step 2: Distinguished unit.} The unramified principal series
\(\mathrm{Ind}_{B(\mathbb{R})}^{GL_{2}(\mathbb{R})}(\chi_{1} \boxtimes \chi_{2})\)
with \(\chi_{1} = \chi_{2} = |\cdot|^{0}\) the trivial character gives
the irreducible spherical unramified principal series; this is the
archimedean analogue of the spherical Iwahori-fixed vector at a
finite prime. Harish-Chandra 1958 ``Spherical functions on a
semisimple Lie group'' identifies this module as the unique irreducible
admissible representation of \(GL_{2}(\mathbb{R})\) with a non-zero
\(K_{\infty}\)-fixed vector at weight \(0\). The unit-comparison map
\(\eta_{\infty}: \mathbb{1} \to u_{\infty}\) is the inclusion of the
unit admissible module (the trivial character of the archimedean
Hecke algebra) into the spherical unramified principal series.

\textbf{Step 3: Langlands--Vogan classification.} Vogan 1981
Theorem 5.1.1 gives the Langlands classification for \(GL_{2}(\mathbb{R})\):
the irreducible admissible \((\mathfrak{g}, K_{\infty})\)-modules are
parametrized by Langlands data, which for \(GL_{2}(\mathbb{R})\)
reduces to the triple \((\lambda, \mu, \varepsilon)\). The classes are:
\begin{itemize}
\item principal series (including the spherical unramified one),
parametrized by \(\lambda \in \mathbb{C}\) and central character
sign \(\varepsilon\);
\item discrete series \(D_{k}\) for \(k \in \mathbb{Z}_{\geq 2}\), with
lowest weight \(k\);
\item limits of discrete series \(D_{1}^{\pm}\);
\item irreducible unitary finite-dimensional modules (only the
trivial module here).
\end{itemize}
Each class is \(K_{\infty}\)-finite in each isotype, hence admissible.
Wallach 1988 Ch.\ 5 organises these into the admissible spectrum.

\textbf{Step 4: Base case identification.} The empty-\(S\) stage of the
restricted tensor product is \(\bigotimes_{v \in \emptyset}\mathcal{C}_{v}
= \mathbb{1}_{\mathrm{Pr}^{\mathrm{L,st}}} = \mathrm{Sp}\). The
transition map \(\iota_{\emptyset \subseteq \{\infty\}}\) inserts the
spherical unit \(u_{\infty}\). The image of \(\mathrm{Sp}\) under
tensoring with \(u_{\infty}\) is \(\mathrm{Sp} \otimes u_{\infty}\), which
inside \(\mathcal{C}_{\infty}\) is the full subcategory generated by
\(u_{\infty}\). Since \(u_{\infty}\) is a \(K_{\infty}\)-finite admissible
module, this subcategory is \(\mathrm{Rep}(K_{\infty})\) as a
subcategory of admissible \((\mathfrak{g}, K_{\infty})\)-modules. The
empty-\(S\) base is therefore \(\mathrm{Rep}(K_{\infty})\), not Vec.
\end{proof}

\subsection{Compatibility with Finite-Place Sph-Fin Structure}

\begin{proposition}[uniformity of the sph-fin definition across places]
\label{v4-arith:lacd:prop:uniformity-sph-fin}
\ClaimStatusProvedHere. The sph-fin condition at an archimedean
place (\Cref{v4-arith:lacd:def:arch-sph-fin}) is the natural
admissibility restriction, and coincides with the finite-prime
sph-fin condition (compact-generated \(\infty\)-category with
finite-type compacts) under the identification of \(K_{\infty}\)-finite
admissible modules with finite-type modules over the archimedean
Hecke algebra.
\end{proposition}

\begin{proof}
Wallach 1988 \S0.A associates to each admissible
\((\mathfrak{g}, K_{\infty})\)-module a module over the archimedean
Hecke algebra \(\mathcal{H}_{\infty}\), which is the convolution
algebra of bi-\(K_{\infty}\)-finite distributions on \(GL_{2}(\mathbb{R})\).
This gives a faithful embedding of admissible
\((\mathfrak{g}, K_{\infty})\)-modules into modules over
\(\mathcal{H}_{\infty}\). Finite-type \(\mathcal{H}_{\infty}\)-modules
correspond to \(K_{\infty}\)-finite admissible modules
(Harish-Chandra 1966 \S4). Hence the archimedean sph-fin
\(\mathcal{C}_{\infty}^{\mathrm{sph\text{-}fin}}\) corresponds to the
full subcategory of compactly generated
\(\mathcal{H}_{\infty}\)-modules, matching the finite-prime definition.
\end{proof}

\begin{corollary}[archimedean factor enters the restricted tensor]
\label{v4-arith:lacd:cor:arch-factor-enters}
\ClaimStatusProvedHere. The archimedean factor \(\mathcal{C}_{\infty}\)
enters the restricted tensor product
\(\bigotimes'_{v \in |\overline{C}|}\mathcal{C}_{v}\) on equal footing
with the finite-prime factors, with distinguished unit
\(u_{\infty}\) as in \Cref{v4-arith:lacd:thm:L0-closed}. The indexing
set is \(V = |\overline{C}| = \mathbb{P} \cup \{\infty\}\), and the
cofinite unramified subset \(V^{\mathrm{sph}}\) includes both
finite unramified primes and the archimedean place.
\end{corollary}

\begin{remark}[archimedean admissibility tightens the \((L_{4})\) extension requirement]
\label{v4-arith:lacd:rem:L4-reconciliation}
The subcondition \((L_{4})\) of \Cref{v4-arith:acd:thm:modular-decomposition}
was stated as ``the archimedean factor extends the finite-place product by
direct tensoring.'' \(\mathrm{L}_{\mathrm{ACD}}.0\) tightens this: the
archimedean factor is Harish-Chandra admissible rather than Iwahori-fixed,
and the empty-\(S\) base case is \(\mathrm{Rep}(K_{\infty})\) rather than
Vec. The original \((L_{4})\) statement is preserved as a corollary;
the tightening is substantive, since no Iwahori subgroup of
\(GL_{2}(\mathbb{R})\) gives the natural finiteness cut at \(v = \infty\).
\end{remark}

\section{\(\mathrm{L}_{\mathrm{ACD}}.1\): Dg-Categorical Colimit Commutation}
\label{v4-arith:lacd:L1}

\subsection{Statement}

\begin{theorem}[\(\mathrm{L}_{\mathrm{ACD}}.1\) closed given F2.c]
\label{v4-arith:lacd:thm:L1-closed}
\ClaimStatusProvedHereConditional \emph{(conditional on F2.c tightened,
\Cref{v4-arith:tate:thm:F2c-main}; F2.c tightened is itself a theorem,
see \Cref{v4-arith:lacd:rem:F2c-closed})}. Let
\(\{(\mathcal{C}_{v}, u_{v}, \eta_{v})\}_{v}\) be a family of pointed
\(E_{2}\)-monoidal stable presentable \(\infty\)-categories (in the
sense of \Cref{v4-arith:sw:def:pointed-En}, with \(\eta_{v}\) an
\(E_{2}\)-algebra map per the tightening of
\Cref{v4-arith:tate:thm:F2c-main})
with \(u_{v} = \mathbb{1}_{\mathcal{C}_{v}}\) at almost every \(v\). Let
\(V^{\mathrm{sph}} \subseteq V\) be a cofinite subset of places.
Assume each \(\mathcal{C}_{v}\) carries a sph-fin subcategory
\(\mathcal{C}_{v}^{\mathrm{sph\text{-}fin}} \subseteq \mathcal{C}_{v}\)
compact-generating, and \(u_{v}\) is compact for \(v \in V^{\mathrm{sph}}\).
Then the restricted tensor product commutes with sph-fin truncation:
\begin{equation}
\label{v4-arith:lacd:eq:colim-commutes}
\bigl(\bigotimes_{v}^{\prime}\mathcal{C}_{v}\bigr)^{\mathrm{sph\text{-}fin}}
\;\simeq\;
\bigotimes_{v}^{\prime}\mathcal{C}_{v}^{\mathrm{sph\text{-}fin}}
\end{equation}
as stable presentable \(\infty\)-categories.
\end{theorem}

\begin{proof}
\textbf{Step 1: Both sides are compactly generated.}
\Cref{v4-arith:tate:thm:F2c-main}(M2) and
\Cref{v4-arith:acd:lem:sph-fin-preservation} together show that
\(\bigotimes'_{v}\mathcal{C}_{v}\) is compactly generated, with compact
subcategory
\(\varinjlim_{S}\bigotimes_{v\in S}\mathcal{C}_{v}^{\omega}\). The
right-hand side of \eqref{v4-arith:lacd:eq:colim-commutes} is the
restricted tensor product of sph-fin categories, itself compactly
generated by \Cref{v4-arith:acd:lem:sph-fin-preservation}.

\textbf{Step 2: Compact objects match.} Lurie HA 5.3.2.9 gives that
a filtered colimit of compactly generated \(\infty\)-categories
along compact-generator-preserving functors is compactly generated,
with compact subcategory the filtered colimit of compact
subcategories. Applied to the defining filtered colimit of the
restricted tensor product:
\begin{equation}
\Bigl(\bigotimes_{v}^{\prime}\mathcal{C}_{v}\Bigr)^{\omega}
\;=\;\varinjlim_{S}\bigotimes_{v\in S}\mathcal{C}_{v}^{\omega}.
\end{equation}

\textbf{Step 3: Sph-fin is the compact subcategory.} By the sph-fin
definition in
\Cref{v4-arith:acd:lem:sph-fin-preservation}, the sph-fin condition
\emph{is} the compact-generation structure enforced by the
distinguished spherical units. Hence
\begin{equation}
\Bigl(\bigotimes_{v}^{\prime}\mathcal{C}_{v}\Bigr)^{\mathrm{sph\text{-}fin}}
\;=\;\mathrm{Ind}\Bigl(\bigl(\bigotimes_{v}^{\prime}\mathcal{C}_{v}\bigr)^{\omega}\Bigr)
\;=\;\mathrm{Ind}\Bigl(\varinjlim_{S}\bigotimes_{v\in S}\mathcal{C}_{v}^{\omega}\Bigr).
\end{equation}

\textbf{Step 4: Ind commutes with filtered colim.} Lurie HTT 5.3.5.15
(Ind-completion preserves filtered colimits of small
\(\infty\)-categories). The Ind of a filtered colimit of small
\(\infty\)-categories is the filtered colimit of Ind-completions.
Hence
\begin{equation}
\mathrm{Ind}\Bigl(\varinjlim_{S}\bigotimes_{v\in S}\mathcal{C}_{v}^{\omega}\Bigr)
\;=\;\varinjlim_{S}\mathrm{Ind}\Bigl(\bigotimes_{v\in S}\mathcal{C}_{v}^{\omega}\Bigr)
\;=\;\varinjlim_{S}\bigotimes_{v\in S}\mathcal{C}_{v}^{\mathrm{sph\text{-}fin}}
\;=\;\bigotimes_{v}^{\prime}\mathcal{C}_{v}^{\mathrm{sph\text{-}fin}}.
\end{equation}

\textbf{Step 5: \(E_{2}\)-monoidal structure matches.} F2.c tightened
(\Cref{v4-arith:tate:thm:F2c-main}(M3)) gives that the
\(E_{2}\)-monoidal structure on the restricted tensor product is
unique up to a contractible space of choices, and the uniqueness is
cofinality-invariant. The \(E_{2}\)-structure on the LHS of
\eqref{v4-arith:lacd:eq:colim-commutes} is inherited from
\(\bigotimes'_{v}\mathcal{C}_{v}\) via restriction to the sph-fin
subcategory; the \(E_{2}\)-structure on the RHS is inherited from
\(\bigotimes'_{v}\mathcal{C}_{v}^{\mathrm{sph\text{-}fin}}\). Both
agree by the universal property of the canonical
\(E_{2}\)-monoidal structure.

Assembling Steps 1--5: both sides of
\eqref{v4-arith:lacd:eq:colim-commutes} are compactly generated
stable presentable \(\infty\)-categories with isomorphic compact
subcategories and isomorphic \(E_{2}\)-monoidal structures. The
isomorphism follows.
\end{proof}

\begin{remark}[role of F2.c tightening]
\label{v4-arith:lacd:rem:F2c-role}
\(\mathrm{L}_{\mathrm{ACD}}.1\) is a consequence of F2.c tightened plus
Lurie HA 5.3.2.9 and HTT 5.3.5.15. F2.c tightened supplies the
universal property of the colimit's \(E_{2}\)-monoidal structure;
HA 5.3.2.9 supplies the sph-fin preservation under the filtered
colimit; HTT 5.3.5.15 supplies the Ind-filtered-colimit exchange. The
three ingredients together close \(\mathrm{L}_{\mathrm{ACD}}.1\). The
source list is entirely Lurie HA + HTT, disjoint from
Fargues--Scholze and from Raskin 2021 JAMS.
\end{remark}

\subsection{Explicit Commutation Square}

\begin{corollary}[commutation square]
\label{v4-arith:lacd:cor:commutation-square}
\ClaimStatusProvedHere. The following square in
\(\mathrm{Pr}^{\mathrm{L,st}}\) commutes up to equivalence:
\begin{equation}
\begin{tikzcd}[column sep=large]
\bigl\{\mathcal{C}_{v}^{\mathrm{sph\text{-}fin}}\bigr\}_{v} \ar[r, "\bigotimes'"] \ar[d, hook] & \bigotimes'_{v}\mathcal{C}_{v}^{\mathrm{sph\text{-}fin}} \ar[d, "\simeq"] \\
\bigl\{\mathcal{C}_{v}\bigr\}_{v} \ar[r, "\bigotimes'"] & \bigl(\bigotimes'_{v}\mathcal{C}_{v}\bigr)^{\mathrm{sph\text{-}fin}}
\end{tikzcd}
\end{equation}
where the vertical arrows are sph-fin inclusion / sph-fin truncation
and the horizontal arrows are the Tate restricted tensor product.
\end{corollary}

\begin{proof}
Immediate from \Cref{v4-arith:lacd:thm:L1-closed} plus functoriality
of the restricted tensor product under compact-generator-preserving
functors.
\end{proof}

\section{\(\mathrm{L}_{\mathrm{ACD}}.2\): Cross-Prime Hecke Coherence}
\label{v4-arith:lacd:L2}

\subsection{Statement}

\begin{theorem}[formal tensor-factor Hecke commutativity]
\label{v4-arith:lacd:thm:L2-closed}
\ClaimStatusProvedHere. For a family
\(\{(\mathcal{C}_{v}, u_{v}, \eta_{v})\}_{v}\) of pointed
\(E_{2}\)-monoidal stable presentable \(\infty\)-categories as in
\Cref{v4-arith:lacd:thm:L1-closed}, equipped with local Hecke
operators \(T_{v} \in \mathrm{End}_{E_{2}}(\mathcal{C}_{v})\) (i.e.\
\(E_{2}\)-algebra endomorphisms of each factor) and archimedean
\(U(\mathfrak{gl}_{2})^{K_{\infty}}\)-action on \(\mathcal{C}_{\infty}\):
\begin{enumerate}
\item For any two distinct places \(v \neq w\), the extensions
\(\widetilde{T}_{v}, \widetilde{T}_{w}: \bigotimes'_{v}\mathcal{C}_{v}
\to \bigotimes'_{v}\mathcal{C}_{v}\) of the local operators by
identity on all other factors commute up to canonical homotopy:
\begin{equation}
\label{v4-arith:lacd:eq:T-commute}
\widetilde{T}_{v} \circ \widetilde{T}_{w}
\;\simeq\;
\widetilde{T}_{w} \circ \widetilde{T}_{v}.
\end{equation}
\item Associativity and unit coherences hold: for any three distinct
places \(v_{1}, v_{2}, v_{3}\), the square associators and unitors
satisfy the standard symmetric-monoidal MacLane pentagon / triangle
coherences (Lurie HA 5.1.2.6).
\item The archimedean
\(U(\mathfrak{gl}_{2})^{K_{\infty}}\)-equivariance commutes with every
finite-prime \(T_{p}\).
\end{enumerate}
\end{theorem}

\begin{remark}[domain: formal versus global Hecke coherence]
\label{v4-arith:lacd:rem:L2-formal-domain}
The theorem establishes commutativity of local endofunctors on the
restricted tensor product and the archimedean equivariance. Global
Hecke functoriality of the automorphic category and the local-global
Whittaker--Bernstein compatibility are distinct properties; both remain
part of (BC-diamond-glue) and are outside the domain of
\Cref{v4-arith:lacd:thm:L2-closed}.
\end{remark}

\begin{proof}
\textbf{Step 1: Tensor-factor disjointness.} Each Hecke operator
\(T_{v}\) acts on the single factor \(\mathcal{C}_{v}\) of the restricted
tensor product. The extension by identity on other factors,
\(\widetilde{T}_{v} := T_{v} \otimes \mathrm{id}_{\neq v}\), is an
endofunctor of \(\bigotimes'_{v}\mathcal{C}_{v}\) that factors
through the tensor factor at \(v\). For distinct places \(v \neq w\), the
endofunctors \(\widetilde{T}_{v}\) and \(\widetilde{T}_{w}\) factor
through disjoint tensor factors.

\textbf{Step 2: Disjoint tensor factors commute up to canonical
homotopy.} Lurie HA 4.8.1 states that endofunctors of a tensor
product \(\mathcal{C} \otimes \mathcal{D}\) that factor through
disjoint tensor factors commute up to a canonical homotopy, because
they correspond to bi-colimit-preserving functors
\(\mathcal{C} \times \mathcal{D} \to \mathcal{C} \otimes \mathcal{D}\)
and the product in the source is symmetric. Applied to
\(\widetilde{T}_{v}\) and \(\widetilde{T}_{w}\) acting on
\(\mathcal{C}_{v} \otimes \mathcal{C}_{w} \otimes
\bigotimes'_{u \neq v, w}\mathcal{C}_{u}\): commutativity holds up to
canonical homotopy.

\textbf{Step 3: Higher coherences.} Lurie HA 5.1.2.6 asserts that
\(E_{2}\)-monoidal stable presentable \(\infty\)-categories form a
\((\infty, 3)\)-category in which the coherence data for a
symmetric-monoidal pattern is contractible. The three-fold
composition of Hecke operators at distinct places \(v_{1}, v_{2}, v_{3}\)
admits a unique coherent associator by the universal property of the
restricted tensor product. The pentagon and triangle coherences are
automatic.

\textbf{Step 4: Archimedean equivariance.} The
\(U(\mathfrak{gl}_{2})^{K_{\infty}}\)-action on \(\mathcal{C}_{\infty}\)
is the central Harish-Chandra action on the archimedean
\((\mathfrak{g},K_{\infty})\)-category. It factors through
\(\mathcal{C}_{\infty}\) alone. For any finite prime \(p\), the Hecke
operator \(T_{p}\) factors through \(\mathcal{C}_{p}\). The two factor
through disjoint tensor factors (\(\mathcal{C}_{\infty}\) and
\(\mathcal{C}_{p}\)). By Step 2, they commute up to canonical
homotopy.

\textbf{Full coherence.} The full coherence of the
symmetric-monoidal pattern is encoded by the requirement that
\(\bigotimes'_{v}\mathcal{C}_{v}\) is an \(E_{2}\)-algebra object in
\(\mathrm{Pr}^{\mathrm{L,st}}\) carrying \(E_{2}\)-algebra endomorphisms
\(\widetilde{T}_{v}\) that factor through disjoint tensor factors. Lurie
HA 5.1.2.4 (Dunn additivity \(E_{2} \otimes E_{2} \simeq E_{4}\)) would
suggest that the composite space of such endomorphisms is naturally
\(E_{4}\)-structured; however, restriction to the subspace of
endomorphisms factoring through individual tensor factors (i.e.\
local Hecke operators) reduces the structure back to \(E_{2}\) per
factor. This is the precise sense in which cross-prime Hecke
coherence holds.
\end{proof}

\subsection{Bernstein--Zelevinsky Compatibility}

\begin{conjecture}[\(\infty\)-categorical Bernstein lift]
\label{v4-arith:lacd:conj:BZ-infty-lift}
\ClaimStatusConjectured. The Bernstein--Zelevinsky decomposition of the
abelian category of smooth \(GL_{2}(\mathbb{Q}_{p})\)-representations
lifts to the stable presentable local Hecke category
\(\mathcal C_p\), with the Bernstein centre acting by orthogonal
idempotent projectors and with the spherical finite subcategory
generated by compact objects finite over the spherical Bernstein
centre.
\end{conjecture}

\begin{proposition}[Bernstein--Zelevinsky local compatibility]
\label{v4-arith:lacd:prop:BZ-compat}
\ClaimStatusProvedHereConditional, conditional on
\Cref{v4-arith:lacd:conj:BZ-infty-lift}. At each finite prime \(p\), the local Hecke
category \(\mathcal{C}_{p}\) decomposes along Bernstein--Zelevinsky
components (Bernstein--Zelevinsky 1976 Ann.\ ENS \S2): each
Bernstein component corresponds to a cuspidal datum
\([M, \sigma]\) of a Levi \(M \subseteq GL_{2}(\mathbb{Q}_{p})\) and
a cuspidal supercuspidal representation \(\sigma\) of \(M\). The Hecke
operator \(T_{p}\) preserves each Bernstein component.
\end{proposition}

\begin{proof}
Bernstein--Zelevinsky 1976 proves the decomposition of the category
of smooth representations of \(GL_{2}(\mathbb{Q}_{p})\) into
Bernstein components. Each component is invariant under the
centre of the Hecke algebra, which includes the spherical Hecke
operator \(T_{p}\) acting on the spherical principal series
component (the unramified component). The passage from this abelian
decomposition to \(\mathcal C_p\) is exactly
\Cref{v4-arith:lacd:conj:BZ-infty-lift}. Under that datum, the
idempotent projectors commute with \(T_p\), so \(T_p\) preserves every
lifted Bernstein component.
\end{proof}

\begin{remark}[role of Bernstein--Zelevinsky]
\label{v4-arith:lacd:rem:BZ-role}
Bernstein--Zelevinsky 1976 supplies the local decomposition of the
smooth representation category. After
\Cref{v4-arith:lacd:conj:BZ-infty-lift}, the same idempotents act on
the stable local Hecke category and are compatible with cross-prime
Hecke coherence, because Hecke operators at distinct primes act on
disjoint tensor factors. The global statement is therefore conditional
on the lift datum; it is not a consequence of the abelian
Bernstein--Zelevinsky theorem alone.
\end{remark}

\section{\(\mathrm{L}_{\mathrm{ACD}}.3\): Sph-Finite Preservation}
\label{v4-arith:lacd:L3}

\subsection{Statement}

\begin{theorem}[\(\mathrm{L}_{\mathrm{ACD}}.3\) automatic given \(\mathrm{L}_{\mathrm{ACD}}.1\)]
\label{v4-arith:lacd:thm:L3-closed}
\ClaimStatusProvedHere. The sph-fin subcategory
\(\mathrm{D}^{\mathrm{sph\text{-}fin}}([GL_{2}(\mathbb{Q}_{v})])\)
is preserved under the restricted-tensor colimit: the colimit of
sph-fin categories embeds in the sph-fin subcategory of the colimit,
and this embedding is an equivalence:
\begin{equation}
\varinjlim_{S}\bigotimes_{v \in S}\mathcal{C}_{v}^{\mathrm{sph\text{-}fin}}
\;\xrightarrow{\sim}\;
\Bigl(\bigotimes_{v}^{\prime}\mathcal{C}_{v}\Bigr)^{\mathrm{sph\text{-}fin}}.
\end{equation}
\end{theorem}

\begin{proof}
Immediate from \Cref{v4-arith:lacd:thm:L1-closed}: the right-hand
side equals \(\bigotimes'_{v}\mathcal{C}_{v}^{\mathrm{sph\text{-}fin}}\),
which by definition of the restricted tensor product is the
filtered colimit of finite tensor products of sph-fin factors,
i.e.\ the left-hand side.
\end{proof}

\begin{corollary}[sph-fin preserved fully]
\label{v4-arith:lacd:cor:sph-fin-preserved}
\ClaimStatusProvedHere. The forgetful functor
\(\mathrm{Pr}^{\mathrm{L,st,sph\text{-}fin}} \to \mathrm{Pr}^{\mathrm{L,st}}\)
preserves the restricted tensor product of pointed \(E_{2}\)-monoidal
categories satisfying
Setup~\ref{v4-arith:tate:setup:data}.
\end{corollary}

\begin{proof}
Consequence of \Cref{v4-arith:lacd:thm:L3-closed}.
\end{proof}

\subsection{Uniformity of the Sph-Fin Subcategory}

\begin{proposition}[uniformity across finite primes]
\label{v4-arith:lacd:prop:uniformity-finite}
\ClaimStatusProvedHereConditional, conditional on
\Cref{v4-arith:lacd:conj:BZ-infty-lift}. The sph-fin subcategory
\(\mathrm{D}^{\mathrm{sph\text{-}fin}}([GL_{2}(\mathbb{Q}_{p})])\) at a
finite prime \(p\) is the full subcategory of \(\infty\)-categorically
compact objects generated by the spherical vector
\(|0\rangle_{p}\), and the compact generators are finite over the
local spherical Hecke algebra.
\end{proposition}

\begin{proof}
\Cref{v4-arith:lacd:conj:BZ-infty-lift} supplies the compact
generation and the finiteness over the spherical Bernstein centre. The
spherical vector generates the unramified component, and
\Cref{v4-arith:lacd:prop:BZ-compat} identifies that component inside
the local Hecke category.
\end{proof}

\begin{proposition}[uniformity at infinity]
\label{v4-arith:lacd:prop:uniformity-infinity}
\ClaimStatusProvedHere. The archimedean sph-fin subcategory
\(\mathrm{D}^{\mathrm{sph\text{-}fin}}((\mathfrak{g}, K_{\infty})\text{-mod}^{\mathrm{adm}})\)
is the full subcategory of compact objects generated by the
spherical unramified principal series \(u_{\infty}\), and the compact
generators are \(K_{\infty}\)-finite admissible modules (equivalently,
finite-type over the archimedean Hecke algebra \(\mathcal{H}_{\infty}\)).
\end{proposition}

\begin{proof}
\Cref{v4-arith:lacd:prop:uniformity-sph-fin} supplies the
identification of \(K_{\infty}\)-finite admissible modules with
finite-type \(\mathcal{H}_{\infty}\)-modules. Wallach 1988 Theorem
0.A.2 gives the Grothendieck-category structure on admissible
modules, hence compact generation of the derived category.
\end{proof}

\section{Consequence: Conditional Reduction of \(\mathrm{L}_{\mathrm{ACD}}\)}
\label{v4-arith:lacd:consequence}

\subsection{The Residual Lemma}

\begin{corollary}[\(\mathrm{L}_{\mathrm{ACD}}\) conditional obstruction decomposition]
\label{v4-arith:lacd:cor:LACD-residual}
\ClaimStatusProvedHereConditional. Assume F2.c tightened
(\Cref{v4-arith:tate:thm:F2c-main}), the
\(\infty\)-categorical Bernstein lift
\(\Cref{v4-arith:lacd:conj:BZ-infty-lift}\), the local categorical
comparison datum for the finite-prime categories, the archimedean
boundary datum of \(\Cref{v4-arith:acd:thm:L4-closed}\), and the
global identification of the restricted tensor product with the
automorphic category. Then the remaining Bernstein-centre part of
\(\mathrm{L}_{\mathrm{ACD}}\) is (BC-diamond-glue). The latter
decomposes by
\Cref{v4-arith:bcglue-k5:thm:decomposition} into (BCglue.local)
(Conditional, requires the unwritten local categorical
comparison (F1.a) via Dotto--Emerton--Gee arXiv:2603.26887),
(BCglue.glue) (Conjectural, requires prismatic faithfully-flat
descent of the Nygaard-convention-aware transition maps), and
(BC-diamond-glue.core) (Conjectural, conditional on global arithmetic
Arinkin--Gaitsgory). Without these inputs, (BC-diamond-glue) alone
does not imply \(\mathrm{L}_{\mathrm{ACD}}\).
\end{corollary}

\begin{proof}
By the initial arithmetic comparison split of \(\mathrm{L}_{\mathrm{ACD}}\)
(\Cref{v4-arith:sw:thm:LACD-split}) plus the residual splitting refinement
(\Cref{v4-arith:sw:wave2}), one has a normalisation implication from
the following data:
\[
\mathrm{L}_{\mathrm{ACD}}.0,\,
\mathrm{L}_{\mathrm{ACD}}.1,\,
\mathrm{L}_{\mathrm{ACD}}.2,\,
\mathrm{L}_{\mathrm{ACD}}.3,\,
\mathrm{L}_{\mathrm{ACD}}.4,\,
\text{local comparison},\,
\text{global identification},\,
\text{(BC-diamond-glue)}
\;\Longrightarrow\;
\mathrm{L}_{\mathrm{ACD}} .
\]
The formal items have the following form:
\begin{itemize}
\item \(\mathrm{L}_{\mathrm{ACD}}.0\): archimedean boundary datum
(\Cref{v4-arith:lacd:thm:L0-closed});
\item \(\mathrm{L}_{\mathrm{ACD}}.1\): theorem given F2.c tightened
(\Cref{v4-arith:lacd:thm:L1-closed});
\item \(\mathrm{L}_{\mathrm{ACD}}.2\): formal Hecke commutativity,
with the finite-prime Bernstein part conditional on
\Cref{v4-arith:lacd:conj:BZ-infty-lift}
(\Cref{v4-arith:lacd:thm:L2-closed});
\item \(\mathrm{L}_{\mathrm{ACD}}.3\): theorem
(\Cref{v4-arith:lacd:thm:L3-closed});
\item \(\mathrm{L}_{\mathrm{ACD}}.4\): theorem (Kneser 1966 and
Platonov--Rapinchuk 1994 for \(SL_{2}\); Tate for the determinant
idele-class quotient).
\end{itemize}
Thus this chapter removes only the formal part of the obstruction.
The categorical comparison data and (BC-diamond-glue) remain separate
inputs.
\end{proof}

\begin{remark}[F2.c tightened is a theorem]
\label{v4-arith:lacd:rem:F2c-closed}
F2.c tightened (\Cref{v4-arith:tate:thm:F2c-main}) is stated as
a \ClaimStatusProvedHere theorem in \Cref{v4-arith:tate-tensor}
(Lurie HA + Fresse 2017, disjoint from Fargues--Scholze). It supplies
only the filtered-colimit and \(E_{2}\)-tensor formalism. It does not
supply the \(\infty\)-categorical Bernstein lift, the local categorical
comparison, the archimedean boundary spectral comparison, the global
identification, or (BC-diamond-glue).
\end{remark}

\subsection{Domain of (BC-diamond-glue)}

\begin{remark}[(BC-diamond-glue) required theorem]
\label{v4-arith:lacd:rem:BC-diamond-glue-domain}
By \Cref{v4-arith:acd:rem:bc-diamond-glue-domain}, (BC-diamond-glue)
has required theorem non-formal theorem: gluing of coherent Springer
sheaves across a product of local Emerton--Gee stacks, plus agreement
with the global Bernstein centre under conjectural global arithmetic
Arinkin--Gaitsgory. This estimate is not an estimate for
\(\mathrm{L}_{\mathrm{ACD}}\) itself. The latter also needs the local
categorical comparison, the \(\infty\)-categorical Bernstein lift, the
archimedean boundary datum, and the global identification.
\end{remark}

\begin{remark}[cumulative reduction]
\label{v4-arith:lacd:rem:cumulative-reduction}
The original \(\mathrm{L}_{\mathrm{ACD}}\) contains the full adelic
descent problem. The modular decomposition of
\Cref{v4-arith:acd:thm:modular-decomposition} isolates a
non-formal theorem Bernstein-centre gluing problem, but it does not
prove that this problem is the only obstruction. After the present
chapter the visible residuals are the local categorical comparison,
the \(\infty\)-categorical Bernstein lift, the archimedean boundary
comparison, the global identification, and (BC-diamond-glue).
\end{remark}

\section{Compatibility with \(\Lambda_{\zeta}(s) = \Gamma_{\mathbb{R}}(s)\zeta(s)\) Factorization}
\label{v4-arith:lacd:xi-compat}

\subsection{The Completed \(L\)-Factor}

The completed Riemann \(\xi\)-function is
\begin{equation}
\label{v4-arith:lacd:eq:xi-def}
\xi(s) \;=\; \tfrac{1}{2}s(s-1)\pi^{-s/2}\Gamma(s/2)\zeta(s) \;=\; \tfrac{1}{2}s(s-1)\cdot\Gamma_{\mathbb{R}}(s)\zeta(s)
\end{equation}
with archimedean local \(L\)-factor
\(\Gamma_{\mathbb{R}}(s) := \pi^{-s/2}\Gamma(s/2)\) and
\begin{equation}
\zeta(s) \;=\; \prod_{p}\frac{1}{1 - p^{-s}} \;=\; \prod_{p}L_{p}(s)
\end{equation}
with finite-prime local \(L\)-factors
\(L_{p}(s) = (1 - p^{-s})^{-1}\). The completed \(\xi\) satisfies the
functional equation
\begin{equation}
\label{v4-arith:lacd:eq:fe}
\xi(s) \;=\; \xi(1 - s).
\end{equation}

Any honest \(\mathrm{L}_{\mathrm{ACD}}\) for \(GL_{2}/\mathbb{Q}\) must
produce a trace class whose Tate comparison sees
\(\Lambda_{\zeta}(s)=\Gamma_{\mathbb R}(s)\zeta(s)\), not merely
\(\zeta(s)\). Multiplication by \(\tfrac12 s(s-1)\) then gives the
entire function \(\xi(s)\).

\subsection{Tate Character Identification}

\begin{proposition}[archimedean character \(=\Gamma_{\mathbb{R}}\)]
\label{v4-arith:lacd:prop:arch-character}
\ClaimStatusProvedHere. The Tate local zeta integral of the
archimedean Gaussian, after the standard factor \(2\) normalization, is
\(\Gamma_{\mathbb{R}}(s) = \pi^{-s/2}\Gamma(s/2)\), the archimedean
local \(L\)-factor of the trivial representation.
\end{proposition}

\begin{proof}
Tate 1950 \S4.3 (Mellin transform of the Gaussian) computes
\begin{equation}
\int_{0}^{\infty}e^{-\pi x^{2}}x^{s}\frac{dx}{x} \;=\; \tfrac{1}{2}\pi^{-s/2}\Gamma(s/2) \;=\; \tfrac{1}{2}\Gamma_{\mathbb{R}}(s).
\end{equation}
The archimedean spherical unramified principal series has Whittaker
model whose \(K_{\infty}\)-fixed vector is the Gaussian (up to
normalisation; Jacquet--Langlands SLN 114 \S1.5). The local
\(L\)-factor at infinity is the Mellin transform of the Gaussian,
hence \(\Gamma_{\mathbb{R}}(s)\).
\end{proof}

\begin{proposition}[finite-prime Tate character \(= L_{p}(s)\)]
\label{v4-arith:lacd:prop:finite-character}
\ClaimStatusProvedHere. With multiplicative Haar measure normalized by
\(\operatorname{vol}(\mathbb Z_p^\times)=1\), the Tate local zeta
integral of the characteristic function of \(\mathbb Z_p\) is
\(L_{p}(s) = (1 - p^{-s})^{-1}\), the finite-prime local \(L\)-factor
of the trivial representation.
\end{proposition}

\begin{proof}
Tate 1950 \S2 computes the local \(p\)-adic integral: the spherical
vector \(|0\rangle_{p}\) evaluates to \((1 - p^{-s})^{-1}\) under the
standard local Tate integral.
\end{proof}

\begin{theorem}[Tate product gives \(\Lambda_{\zeta}(s)\)]
\label{v4-arith:lacd:thm:xi-reproduced}
\ClaimStatusProvedHere. The product of the normalized Tate local
characters over all places of \(\mathbb Q\) is
\begin{equation}
\prod_v Z_v(s)
\;=\; \Gamma_{\mathbb{R}}(s)\prod_{p}L_{p}(s)
\;=\; \Gamma_{\mathbb{R}}(s)\zeta(s)
\;=\; \frac{\xi(s)}{\tfrac{1}{2}s(s-1)},
\end{equation}
as a meromorphic identity. Thus the Tate comparison target is
\(\Lambda_{\zeta}(s)=\Gamma_{\mathbb{R}}(s)\zeta(s)\).
\end{theorem}

\begin{proof}
This is Tate's product formula for the standard adelic test function:
\begin{equation}
\prod_v Z_v(s)
\;=\;\Gamma_{\mathbb{R}}(s)\prod_{p}L_{p}(s)
\;=\;\Gamma_{\mathbb{R}}(s)\zeta(s).
\end{equation}
The standard completed \(\xi\)-function is obtained from this
meromorphic factor by the elementary multiplier \(\tfrac12 s(s-1)\).
\end{proof}

\subsection{Functional Equation Compatibility}

\begin{corollary}[functional equation preserved]
\label{v4-arith:lacd:cor:fe-preserved}
\ClaimStatusProvedHere. The completed Tate product
\(\Lambda_{\zeta}(s)=\Gamma_{\mathbb{R}}(s)\zeta(s)\) satisfies the
functional equation
\(\Gamma_{\mathbb{R}}(s)\zeta(s) = \Gamma_{\mathbb{R}}(1-s)\zeta(1-s)\),
equivalently \(\xi(s) = \xi(1-s)\).
\end{corollary}

\begin{proof}
\Cref{v4-arith:lacd:thm:xi-reproduced} gives
\(\Lambda_{\zeta}\). Tate 1950 identifies the Fourier transform of the
standard adelic test function with itself and obtains
\(\Lambda_{\zeta}(s)=\Lambda_{\zeta}(1-s)\). Multiplication by
\(\tfrac12 s(s-1)\) gives the displayed equation for \(\xi\).
\end{proof}

\begin{remark}[significance of the parallel archimedean factor]
\label{v4-arith:lacd:rem:arch-significance}
Without the parallel archimedean factor, the restricted tensor
product \(\bigotimes'_{p<\infty}\mathcal{C}_{p}\) could at most
compare with \(\zeta(s)\), not with \(\Lambda_{\zeta}(s)\). The
finite-prime product satisfies no functional equation of its own; only the
completed Tate product
\(\Lambda_{\zeta}(s) = \Gamma_{\mathbb{R}}(s)\zeta(s)\) does.
Hence \(\mathrm{L}_{\mathrm{ACD}}.0\) (archimedean stratum) is a
\emph{mathematical necessity}, not a cosmetic addition: it is
required before any restricted-tensor construction can be compared
with the completed Tate character.
\end{remark}

\begin{remark}[trace comparison is extra structure]
\label{v4-arith:lacd:rem:trace-comparison-extra}
The Tate calculation above is not a computation of Hochschild trace in
the local Hecke categories. A categorical realization must still
construct a trace-to-Tate comparison map and prove that the chosen
local units map to the standard local test functions.
\end{remark}

\section{Residual Conditions}
\label{v4-arith:lacd:residual}

\subsection{Residual LACD Conditions}

\(\mathrm{L}_{\mathrm{ACD}}\) contains the local categorical comparison, the
\(\infty\)-categorical Bernstein lift, the archimedean boundary
comparison, the global identification, and (BC-diamond-glue).

\begin{table}[h]
\centering
\small
\begin{tabular}{l|l|l}
Subcondition & Form & Source \\
\hline
\(\mathrm{L}_{\mathrm{ACD}}.0\) (archimedean stratum) & \textbf{Boundary datum} & Harish-Chandra 1966; Vogan 1981; Wallach 1988 \\
\(\mathrm{L}_{\mathrm{ACD}}.1\) (colimit commutation) & \textbf{Theorem} (F2.c tightened) & Lurie HA 5.3.2.9 + F2.c \\
\(\mathrm{L}_{\mathrm{ACD}}.2\) (cross-prime Hecke) & \textbf{Formal theorem; conditional Bernstein part} & Lurie HA; \(\infty\)-categorical BZ lift \\
\(\mathrm{L}_{\mathrm{ACD}}.3\) (sph-fin preservation) & \textbf{Theorem} (from .1) & Immediate \\
\(\mathrm{L}_{\mathrm{ACD}}.4\) (derived strong approximation and determinant quotient) & \textbf{Theorem} & Kneser 1966; Platonov--Rapinchuk 1994; Tate 1950 \\
F2.c tightened & \textbf{Theorem under pointed \(E_{2}\) hypotheses} (see Rem.~\ref{v4-arith:lacd:rem:F2c-closed}) & Lurie HA; Fresse 2017 \\
\(\infty\)-categorical BZ lift & \textbf{Conjectural} & \Cref{v4-arith:lacd:conj:BZ-infty-lift} \\
local categorical comparison and global identification & \textbf{Open} & F1.a and \(\mathrm{L}_{\mathrm{ACD}}\) input \\
(BC-diamond-glue) & \textbf{Conjectural} & Ben-Zvi--Chen--Helm--Nadler arXiv:2010.02321; construction \(\varepsilon\) \\
\end{tabular}
\caption{Form of \(\mathrm{L}_{\mathrm{ACD}}\) components. Formal
pieces are separated from comparison data.}
\label{v4-arith:lacd:tbl:form}
\end{table}

\subsection{F1 Closure Chain}

\begin{theorem}[F1 residual after local comparison and descent hypotheses]
\label{v4-arith:lacd:thm:F1-residual}
\ClaimStatusNeedsVerification. Assume F2.c tightened
(\Cref{v4-arith:tate:thm:F2c-main}), the local categorical
comparison (F1.a), the \(\infty\)-categorical Bernstein lift, the
archimedean boundary comparison, and the global automorphic descent
hypotheses.
Then F1 (global arithmetic
Arinkin--Gaitsgory for \(GL_{2}/\mathbb{Q}\),
\Cref{v4-arith:conj:F1-global-AG}) reduces to
(BC-diamond-glue) (\Cref{v4-arith:acd:cj:bc-diamond-glue-concrete}).
Without those hypotheses, (BC-diamond-glue) alone does not close F1.
\end{theorem}

\begin{proof}
By \Cref{v4-arith:thm:F1-reduction}, F1 reduces to
(F1.a) + (F1.b) + (F1.c) + (F1.d). The item (F1.a) is an assumed
local categorical comparison; (F1.c) is the Emerton--Gee stack input.
Items (F1.b) and (F1.d) combine into \(\mathrm{L}_{\mathrm{ACD}}\)
(\Cref{v4-arith:cl:cor:F1-single-lemma}). By
\Cref{v4-arith:lacd:cor:LACD-residual}, \(\mathrm{L}_{\mathrm{ACD}}\)
reduces to (BC-diamond-glue) only after the comparison hypotheses
listed in the statement are supplied.
\end{proof}

\begin{remark}[form of F1 after this chapter]
\label{v4-arith:lacd:rem:F1-form}
F1 does not reduce to (BC-diamond-glue) alone. The reduction still
contains the local comparison, the \(\infty\)-categorical Bernstein
lift, the archimedean boundary comparison, and the global descent
input. Construction \(\varepsilon\) addresses only the Bernstein-centre
gluing part.
\end{remark}

\subsection{Hypotheses Inventory}

\begin{table}[h]
\centering
\small
\begin{tabular}{l|l|l}
Hypothesis & Role in \(\mathrm{L}_{\mathrm{ACD}}\) reduction & Source \\
\hline
F2.c tightened & Universal property for colimit \(E_{2}\)-structure & Lurie HA, pointed \(E_{2}\) hypotheses \\
Lurie HA 5.3.2.9 & Sph-fin preservation under filtered colimit & Lurie HA \\
Lurie HA 5.1.2.4 & Dunn additivity for \(E_{n}\)-operads & Lurie HA \\
Lurie HA 5.1.2.26 & Unit centrality in \(E_{n}\)-monoidal cat & Lurie HA \\
Lurie HTT 5.3.5.15 & Ind-completion preserves filtered colim & Lurie HTT \\
Harish-Chandra 1966 & Admissibility of \((\mathfrak{g}, K_{\infty})\)-mod & Primary source \\
Vogan 1981 & Langlands--Vogan classification at \(\mathbb{R}\) & Primary source \\
Wallach 1988/1992 & Real reductive groups spectrum theory & Primary source \\
Jacquet--Langlands SLN 114 & Archimedean local representation & Primary source \\
Bernstein--Zelevinsky 1976 & Bernstein decomposition of \(p\)-adic Hecke & Primary source \\
\(\infty\)-categorical BZ lift & Stable lift of Bernstein projectors & \Cref{v4-arith:lacd:conj:BZ-infty-lift}; conjectural \\
Tate 1950 & Adelic Mellin/theta and functional equation & \(\xi\)-compatibility \\
Kneser 1966 & Strong approximation for \(SL_{2}\) & \(\mathrm{L}_{\mathrm{ACD}}.4\) derived subgroup theorem \\
Tate 1950 & Idele-class determinant quotient for \(GL_{2}\) & \(\mathrm{L}_{\mathrm{ACD}}.4\) determinant line \\
(BC-diamond-glue) & Bernstein-centre gluing & Construction \(\varepsilon\), not yet theorem \\
\end{tabular}
\caption{Hypotheses used in this chapter. The non-theorem inputs also
include the local comparison, the archimedean boundary comparison, and
the global identification.}
\label{v4-arith:lacd:tbl:hypotheses}
\end{table}

\section{Conditional Reduction of \(\mathrm{L}_{\mathrm{ACD}}\)}
\label{v4-arith:lacd:summary}

The modular decomposition (\Cref{v4-arith:acd:thm:modular-decomposition},
\Cref{v4-arith:sw:thm:LACD-split}, \Cref{v4-arith:sw:wave2}) identifies
five subconditions
\(\mathrm{L}_{\mathrm{ACD}}.0\)--\(\mathrm{L}_{\mathrm{ACD}}.4\), with
\(\mathrm{L}_{\mathrm{ACD}}.4\) combines Kneser's strong approximation
for \(SL_{2}\) with Tate's determinant quotient for \(GL_{2}\). The
formal pieces treated here are:
\begin{itemize}
\item \(\mathrm{L}_{\mathrm{ACD}}.0\) as archimedean boundary datum via Harish-Chandra admissibility,
Vogan 1981 classification, and Wallach 1988/1992 real-reductive-groups
spectrum theory (\Cref{v4-arith:lacd:thm:L0-closed});
\item \(\mathrm{L}_{\mathrm{ACD}}.1\) via F2.c tightened, Lurie HA
5.3.2.9, and HTT 5.3.5.15
(\Cref{v4-arith:lacd:thm:L1-closed});
\item \(\mathrm{L}_{\mathrm{ACD}}.2\) (formal tensor-factor commutativity)
via Lurie HA 5.1.2.4 (Dunn additivity), with finite-prime Bernstein
compatibility conditional on \Cref{v4-arith:lacd:conj:BZ-infty-lift}
(\Cref{v4-arith:lacd:thm:L2-closed});
\item \(\mathrm{L}_{\mathrm{ACD}}.3\) from \(\mathrm{L}_{\mathrm{ACD}}.1\)
(\Cref{v4-arith:lacd:thm:L3-closed}).
\end{itemize}
The residual of \(\mathrm{L}_{\mathrm{ACD}}\) is not only
(BC-diamond-glue). It also contains the local categorical comparison,
the \(\infty\)-categorical Bernstein lift, the archimedean boundary
comparison, and the global identification
(\Cref{v4-arith:lacd:cor:LACD-residual}).

The parallel archimedean factor
(\(\mathrm{L}_{\mathrm{ACD}}.0\)) is required: without it, the
restricted tensor product can at most compare with \(\zeta(s)\), not
with \(\Gamma_{\mathbb{R}}(s)\zeta(s)\), and the Tate functional
equation is invisible
(\Cref{v4-arith:lacd:rem:arch-significance}).

F1 (global arithmetic Arinkin--Gaitsgory for \(GL_{2}/\mathbb{Q}\))
reduces to (BC-diamond-glue) only under the stated comparison and
descent hypotheses
(\Cref{v4-arith:lacd:thm:F1-residual}).

The source list for \(\mathrm{L}_{\mathrm{ACD}}.0\),
\(\mathrm{L}_{\mathrm{ACD}}.1\), and \(\mathrm{L}_{\mathrm{ACD}}.3\) is
disjoint from the Fargues--Scholze list: Harish-Chandra,
Jacquet--Langlands, Vogan, Wallach (archimedean representation
theory), Lurie HA, Fresse 2017 (higher operadic machinery), and
Bernstein--Zelevinsky 1976 (Bernstein decomposition). The
\(\infty\)-categorical Bernstein lift is not proved from these sources;
it is the separate conjectural datum
\Cref{v4-arith:lacd:conj:BZ-infty-lift}.
